\section{Introduzione}

\subsection{Descrizione generale del progetto}
Il progetto prevede la progettazione, lo sviluppo e il rilascio di un'applicazione mobile Android destinata a supportare attivamente le operazioni del \textbf{Laboratorio in Campo}, un'attività didattica organizzata per gli studenti del secondo anno del corso di Scienze Ambientali. Il sistema mira a digitalizzare e ottimizzare il processo di raccolta dati (fino ad ora effettuato tramite carta e penna) fornendo uno strumento affidabile per la registrazione, condivisione e successiva elaborazione delle informazioni ambientali. Gli stakeholder principali del progetto includono gli studenti (utenti finali), i docenti (committenti e supervisori) e il team di sviluppo. Per una visione dettagliata del dominio applicativo, dei requisiti utente e del contesto generale, si rimanda al documento D1 - Studio di Fattibilità.
\subsection{Scopo e obiettivi del piano}
% TODO: a cosa serve questo piano, a chi è rivolto e quali obiettivi di progetto (tempi, costi, qualità) intende governare.
Il presente Piano di Progetto ha lo scopo di definire in modo rigoroso le modalità di esecuzione, monitoraggio e controllo del progetto di sviluppo software. Il documento è rivolto principalmente ai membri del team di sviluppo, al Project Manager e ai docenti responsabili della valutazione. 
L'obiettivo pratico del corso è la realizzazione in team di un'app Android a supporto del Laboratorio in Campo per gli studenti di Scienze Ambientali. 
L'app dovrà gestire principalmente le seguenti funzionalità:
\begin{itemize}
    \item Raccolta dati (piante, rocce, dati fisici/chimici).
    \item Condivisione di materiali e comunicazione tra i membri del team.
    \item Utilizzo dei sensori dello smartphone (GPS, fotocamera, altitudine).
\end{itemize}


In particolare, gli utenti sono individuati come docenti e studenti, ognuno con ruoli e permessi differenti. Il risultato atteso è un'applicazione funzionante, testata e documentata, che possa essere utilizzata durante le attività di laboratorio in campo.
Gli studenti sono in grado di:
\begin{itemize}
    \item Unirsi a un gruppo di lavoro.
    \item Inserire dati raccolti sul campo, tramite l'upload di foto, scrivere note testuali e modificare i dati inseriti e visualizzare i dati raccolti da altri membri del gruppo.
    \item Visualizzare i dati raccolti in tempo reale.
    \item Generare report in formato xlsx dei dati raccolti e scrivere una relazione finale.
\end{itemize}


I docenti sono in grado di:
\begin{itemize}
    \item Creare gruppi di lavoro e invitare gli studenti a unirsi.
    \item Visualizzare i dati raccolti dagli studenti in tempo reale.
    \item Generare report in formato xlsx dei dati raccolti.

\end{itemize}

\subsection{Definizioni, acronimi e riferimenti}
% TODO: glossario (WBS, PERT, Gantt, ecc.) e riferimenti ai documenti D0 e D1.

\subsubsection*{Definizioni e Acronimi}
Al fine di sciogliere in partenza possibili ambiguità, di seguito si riportano le definizioni e gli acronimi dei termini tecnici utilizzati nel presente Piano di Progetto:

\begin{itemize}
    \item \textbf{WBS (Work Breakdown Structure)}: Struttura di scomposizione del progetto in attività, funzioni e task elementari.
    \item \textbf{PERT (Program Evaluation and Review Technique)}: Diagramma che rappresenta le attività del progetto e le relazioni tra di esse, aiutando a identificare il percorso critico.
    \item \textbf{Gantt}: Diagramma che mostra le attività del progetto nella scala temporale, permettendo di visualizzare sovrapposizioni, dipendenze e durata.
    \item \textbf{SRS (Software Requirements Specification)}: Documento di specifica dei requisiti software.
    \item \textbf{App}: Applicazione mobile per Android.
    \item \textbf{D0, D1, ecc.}: Nomenclatura utilizzata per identificare i vari Deliverable del progetto.
\end{itemize}

\subsubsection*{Riferimenti}
Di seguito sono elencati i documenti e i materiali a cui ci si riferisce nel presente piano:

\begin{itemize}
    \item \textbf{Documenti interni}:
    \begin{itemize}
        \item \textit{Deliverable D0}: Documento contenente la definizione dei gruppi di lavoro, l'elenco dei partecipanti e i relativi CV.
        \item \textit{Deliverable D1}: Documento contenente l'analisi dei requisiti e il contesto generale del progetto.
    \end{itemize}
    \item \textbf{Documenti esterni e manuali}:
    \begin{itemize}
        \item Materiale didattico del corso di Ingegneria del Software.
        \item Documentazione ufficiale per lo sviluppo Android.
    \end{itemize}
\end{itemize}

\subsection{Struttura del documento}
La Sezione~2 descrive l'organizzazione del progetto. La Sezione~3 riporta i processi gestionali e la Sezione~4 i processi tecnici.
La Sezione~5 contiene la pianificazione delle attività e delle risorse. Le appendici raccolgono il materiale di riferimento.
